\label{chap:83-10}

% Unidad U016. Registro firmado, lectura terminal y reconstruccion de K.
% Redaccion autonoma desde la autoridad material 1063 y su sucesor V2.

\section{Génération initiale de l'état dodécaphasique signé et reconstruction du sceau \(K\)}
\label{p12-83-c10-s1:chap:registro-dodecafasico-hadamard-k}

La construction du sceau entier ne commence ni par \(\alpha\), ni par son inverse
ni par un tableau métrologique. Son domaine est l’état terminal enrichi
du Topological Phase Kernel obtenu après la chaîne déjà construite
\begin{equation}
 \boxed{
 \mathrm{APP}\longrightarrow\mathrm{TRIT}\longrightarrow\mathrm{TPK}
 \longrightarrow\Omega_{\rm seed}
 \xrightarrow{T_{\min}}\mathcal U_6^{\rm cat}
 \xrightarrow{\rho_3}W_{243}
 \longrightarrow w_{30}\longrightarrow R_{36}
 \longrightarrow G_9\longrightarrow x_{\rm term}.}
 \label{p12-83-c10-s1:eq:cadena-upstream-registro-dodecafasico}
\end{equation}
Les flèches précédentes ont été définies avec leurs domaines respectifs :
paires de germes, signatures, émissions, mots ternaires, frontières,
histoires compatibles et holonomie nonadique. Cette unité ne choisit à nouveau
aucune de ces données. Elle construit les cartes entières du même état
terminal et démontre leur équivalence.

\subsection{Traces orientées sur des fenêtres de longueur \(108\)}

Soit
\[
 \Theta_{108}=\mathbb Z/108\mathbb Z
\]
le calendrier observable. Ses douze fenêtres nonadiques sont
\[
 E_m=\{9m+1,\ldots,9m+9\},
 \qquad 0\leq m<12.
\]
La partition est disjointe et exhaustive :
\[
 \{1,\ldots,108\}=\bigsqcup_{m=0}^{11}E_m.
\]
L’ensemble des codes de direction qui publient un événement est fixé comme
\[
 D_{\rm evt}=\{2,4,6,8\}.
\]
L’orientation des douze fenêtres est
\begin{equation}
 \Sigma_{12}=(+,+,+,-,-,-,+,+,+,-,-,-).
 \label{p12-83-c10-s1:eq:signatura-doce-ventanas}
\end{equation}
Si \(d_t\) est la direction codée au pas \(t\), l’événement signé
de la fenêtre \(m\) est
\begin{equation}
 e_m=\Sigma_{12}(m)\,
 \#\{t\in E_m:d_t\in D_{\rm evt}\}.
 \label{p12-83-c10-s1:eq:evento-firmado-ventana}
\end{equation}
La formule distingue le nombre non négatif d’événements et le signe de
l’orientation. Elle ne remplace pas une direction cardinale par un signe : les deux
coordonnées restent présentes dans l’état TPK.

\begin{definition}[Registre enrichi dodécaphasique]
\label{p12-83-c10-s1:def:registro-enriquecido-dodecafasico}
Le registre transporté par la trace est
\[
 \mathcal R_{12}(x)=
 \bigl((E_m,\tau_m,\sigma_m,\kappa_m,\Lambda_m)\bigr)_{m=0}^{11},
\]
où \(\tau_m\) est l’orientation TRIT, \(\sigma_m\) la signature de
frontière, \(\kappa_m\) la retenue et \(\Lambda_m\) le segment de registre généalogique
qui conserve la route. Il existe une chaîne de projections
\[
 \mathcal X_{\rm TPK}^{\rm enr}
 \longrightarrow\mathcal R_{12}
 \longrightarrow\Theta_{108},
\]
mais aucune des deux flèches n’est une identification.
\end{definition}

En particulier, \(\Theta_{108}\) conserve uniquement la position du calendrier.
Le registre \(\mathcal R_{12}\) conserve en outre l’orientation, la signature,
la retenue et la provenance. L’état complet conserve encore le cylindre,
la frontière, le quotient, le terminal et la mémoire verticale.

\subsection{Carte entière de l’opposition dodécaphasique}

Pour \(0\leq i<6\), les positions opposées déterminent
\begin{equation}
 p_i=e_i+e_{i+6},
 \qquad
 a_i=e_i-e_{i+6}.
 \label{p12-83-c10-s1:eq:pares-opuestos-registro}
\end{equation}
La charge centrale et les cinq marges relatives sont
\begin{equation}
 s_C=\sum_{i=0}^{5}p_i,
 \qquad
 M_i=p_i-p_5\quad(0\leq i<5).
 \label{p12-83-c10-s1:eq:carga-margenes-registro}
\end{equation}

\begin{definition}[Carte intégrale \(1+5+6\)]
\label{p12-83-c10-s1:def:carta-integral-156}
On définit
\begin{equation}
 \mathcal L_{12}(e_0,\ldots,e_{11})
 =\bigl(s_C,M_0,\ldots,M_4,a_0,\ldots,a_5\bigr).
 \label{p12-83-c10-s1:eq:carta-integral-156}
\end{equation}
La partition des coordonnées est
\[
 1+5+6=12.
\]
\end{definition}

\begin{theorem}[Inversion exacte de la carte entière]
\label{p12-83-c10-s1:thm:inversion-carta-integral-156}
Un tuple
\[
 \ell=(s_C,M_0,\ldots,M_4,a_0,\ldots,a_5)\in\mathbb Z^{12}
\]
appartient à l’image de \(\mathcal L_{12}\) si et seulement si
\begin{align}
 s_C-\sum_{i=0}^{4}M_i&\equiv0\pmod6,
 \label{p12-83-c10-s1:eq:congruencia-seis-carta}\\
 p_i&\equiv a_i\pmod2
 \quad(0\leq i<6),
 \label{p12-83-c10-s1:eq:congruencia-paridad-carta}
\end{align}
où
\begin{align}
 p_5&=\frac{s_C-\sum_{i=0}^{4}M_i}{6},
 \label{p12-83-c10-s1:eq:inversa-p5-carta}\\
 p_i&=p_5+M_i\quad(0\leq i<5).
 \label{p12-83-c10-s1:eq:inversa-pi-carta}
\end{align}
Dans ce cas, l’unique antécédent entier est
\begin{equation}
 e_i=\frac{p_i+a_i}{2},
 \qquad
 e_{i+6}=\frac{p_i-a_i}{2}.
 \label{p12-83-c10-s1:eq:inversa-eventos-carta}
\end{equation}
\end{theorem}

\begin{proof}
En sommant \(M_i=p_i-p_5\) pour \(0\leq i<5\), nous obtenons
\[
 \sum_{i=0}^{4}M_i
 =\sum_{i=0}^{4}p_i-5p_5
 =s_C-6p_5,
\]
ce qui impose \eqref{p12-83-c10-s1:eq:inversa-p5-carta} et la congruence modulo six.
Une fois les six \(p_i\) reconstruits, le système
\[
 p_i=e_i+e_{i+6},
 \qquad a_i=e_i-e_{i+6}
\]
admet l’unique solution \eqref{p12-83-c10-s1:eq:inversa-eventos-carta}. Cette solution est
entière exactement lorsque \(p_i\) et \(a_i\) ont la même parité.
Les deux familles de congruences sont donc nécessaires et suffisantes.
\end{proof}

Ce théorème est antérieur à la clôture numérique de \(\alpha\). Il explique comment la
trace orientée produit une carte entière réversible, et pourquoi il ne suffit pas
de conserver le calendrier réduit.

\subsection{Sélection de l’état terminal commun}

Les trois sections compatibles atteignent la profondeur terminale avec des catalogues
de cardinalités
\[
 |\mathcal T_\pi|=196,
 \qquad
 |\mathcal T_e|=197,
 \qquad
 |\mathcal T_\varphi|=196.
\]
Le produit contient
\begin{equation}
 196\cdot197\cdot196=7\,567\,952
 \label{p12-83-c10-s1:eq:cardinal-producto-terminal}
\end{equation}
triplets ordonnés.

La marge de colonnes du lecteur terminal est
\[
 q_\partial=(1,2,2,1,2,0)\in\mathbb F_3^6,
\]
et l’orientation des lignes du canal conjoint est
\[
 \sigma=(1,1,-1)=(1,1,2)\in\mathbb F_3^3.
\]
La première donnée agit sur les colonnes ; la seconde sur les lignes. L’une ne se déduit pas
de l’autre.

Les signatures locales de Gram, de norme, de poids et de distance réduisent
\eqref{p12-83-c10-s1:eq:cardinal-producto-terminal} à \(331\) triplets. La marge
\(q_\partial\) laisse deux indices :
\[
 (120,197,157),
 \qquad
 (129,197,148).
\]
Leurs matrices sont
\[
 B_0=
 \begin{pmatrix}
 0&2&0&2&2&0\\
 0&0&1&2&2&0\\
 1&0&1&0&1&0
 \end{pmatrix},
 \qquad
 B_1=
 \begin{pmatrix}
 0&2&1&0&2&0\\
 0&0&1&2&2&0\\
 1&0&0&2&1&0
 \end{pmatrix}.
\]
Les sommes des lignes sont
\[
 (0,2,0),
 \qquad
 (2,2,1).
\]
Comme
\[
 \langle\sigma\rangle
 =\{(0,0,0),(1,1,2),(2,2,1)\},
\]
seul \(B_1\) satisfait la condition axiale.

\begin{theorem}[Unicité de la publication terminale visible]
\label{p12-83-c10-s1:thm:unicidad-publicacion-terminal-visible}
La sélection conjointe détermine
\begin{equation}
 B_{\rm term}=
 \begin{pmatrix}
 021020\\001220\\100210
 \end{pmatrix},
 \label{p12-83-c10-s1:eq:matriz-terminal-visible}
\end{equation}
correspondant au triplet \((129,197,148)\). Aucune évaluation décimale de
\(\alpha\) n’intervient dans le recensement ni dans les deux conditions de sélection.
\end{theorem}

\begin{proof}
Le recensement et la marge laissent exactement \(B_0,B_1\). La charge de \(B_0\)
n’appartient pas à \(\langle\sigma\rangle\), tandis que celle de \(B_1\) est
\(2\sigma\). La seconde candidate est donc la seule qui satisfasse
simultanément les deux exigences typées.
\end{proof}

\subsection{Publications archimédienne et exceptionnelle de l'état commun : codomaines non équivalents et généalogie partagée}

Soit
\[
 x_{\rm term}\in
 \mathcal X_{\rm TPK}^{\rm term,enr}
 \subseteq\mathcal X_{\rm TPK}^{[108]}.
\]
La lecture visible est
\[
 \pi_{\rm term}(x_{\rm term})=B_{\rm term}.
\]
La carte signée conserve une autre coordonnée du même état.

\begin{definition}[Lecture causale signée terminale]
\label{p12-83-c10-s1:def:subespacio-firmado-terminal}
Le domaine n’incorpore comme coordonnées d’entrée ni \(K\) ni \(U\).
L’état \(x_{\rm term}\), produit par la composition complète
\(\mathrm{APP}\to\mathrm{TRIT}\to\mathrm{TPK}\), contient le registre
terminal enrichi, et l’extracteur de ses douze fenêtres orientées est
\begin{equation}
 R_{12}:\mathcal X_{\rm TPK}^{\rm term,enr}
 \longrightarrow\mathcal R_{12}^{\rm or},
 \qquad
 R_{12}(x_{\rm term})
 =\bigl((E_m,\tau_m,\sigma_m,\kappa_m,\Lambda_m)\bigr)_{m=0}^{11}.
 \label{p12-83-c10-s1:eq:subespacio-firmado-terminal}
\end{equation}
Le lecteur de canaux \(90/120\) du registre émet, sans recevoir le résultat
comme donnée,
\[
 E_{108}^{90,120}:\mathcal R_{12}^{\rm or}
 \longrightarrow\mathbb Z^{12}\times\mathbb Z^{12}\times\mathbb Z,
 \qquad
 E_{108}^{90,120}\bigl(R_{12}(x_{\rm term})\bigr)
 =(b_{90},b_{120},Q_{\rm TPK}),
\]
avec
\begin{align*}
 b_{90}={}&(-495,-116,-684,108,601,-90,
             -173,-88,313,560,-397,461),\\
 b_{120}={}&(-425,-281,-481,671,-255,30,
              475,-543,680,251,6,-128),\\
 Q_{\rm TPK}={}&6263.
\end{align*}
Pour \(r=1,2,3\) et \(j=0,1,2,3\), on définit
\[
\begin{aligned}
 B_r&=\sum_{j=0}^{3}b_{120,r+3j},
 &d_{r,j}&=b_{90,r+3j},\\
 A_1&=\frac{Q_{\rm TPK}+2B_1+B_2}{3},
 &A_2&=A_1-B_1,
 &A_3&=A_2-B_2.
\end{aligned}
\]
La publication \(\Pi_H\) forme les trois blocs
\[
 U^{(r)}=
 \bigl(A_r,d_{r,1}-d_{r,3},d_{r,0}+d_{r,2},
 d_{r,0}-d_{r,2}\bigr),
\]
et les réentrelace dans l’ordre dodécaphasique, c’est-à-dire
\[
 \Pi_H(b_{90},b_{120},Q_{\rm TPK})
 :=\operatorname{reint}_3\bigl(U^{(1)},U^{(2)},U^{(3)}\bigr)=U.
\]
Soit \(\mathcal H_{12}\) la
transformation entière qui sera construite dans la section suivante. Le module
de publication est
\begin{equation}
 \mathcal U_{12}^{\rm sgn}
 :=\operatorname{im}
 \bigl(\mathcal H_{12}:\mathbb Z^{12}\to\mathbb Z^{12}\bigr).
 \label{p12-83-c10-s1:eq:modulo-publicacion-firmada}
\end{equation}
La lecture signée est, causalement, la composition
\begin{equation}
 R_{\rm sgn}:=
 \Pi_H\circ E_{108}^{90,120}\circ R_{12}:
 \mathcal X_{\rm TPK}^{\rm term,enr}
 \longrightarrow\mathcal U_{12}^{\rm sgn},
 \qquad
 R_{\rm sgn}(x_{\rm term})=U.
 \label{p12-83-c10-s1:eq:lector-firmado-terminal}
\end{equation}
Ce n’est qu’après cette sortie que l’on reconstruit
\(K=\mathcal H_{12}^{-1}U\) ; l’identité
\(U=\mathcal H_{12}K\) certifie cette reconstruction et ne définit pas le domaine
du lecteur.
\end{definition}

L’écriture conjointe correcte est
\begin{equation}
 x_{\rm term}
 \xrightarrow{(\pi_{\rm term},R_{\rm sgn})}
 (B_{\rm term},U_{12}^{\rm sgn}).
 \label{p12-83-c10-s1:eq:doble-publicacion-terminal}
\end{equation}
On n’interpose pas une flèche
\(B_{\rm term}\to U_{12}^{\rm sgn}\) : la matrice visible a précisément oublié
les coordonnées d’orientation, d’opposition, de quotient, de retenue et de
registre généalogique utilisées par la carte signée. La source des deux publications est
l’état enrichi produit par \eqref{p12-83-c10-s1:eq:cadena-upstream-registro-dodecafasico}.

L’évaluation du registre signé donne
\begin{equation}
 \boxed{
 U=(2378,1406,2479,-452,998,-551,
 -668,-204,-371,-322,-28,-997).}
 \label{p12-83-c10-s1:eq:vector-u-firmado-autonomo}
\end{equation}

\subsection{Naturalité par rapport à la profondeur}

Pour \(N'\geq N\geq108\), les troncatures du système signé sont
\begin{equation}
 \tau_{N',N}^{\rm sgn}
 x_{\rm term}^{[N']}
 =\tau_{N',N}x_{\rm term}^{[N']}.
 \label{p12-83-c10-s1:eq:truncamiento-preserva-ku-autonomo}
\end{equation}
Par conséquent, la troncature agit uniquement sur l’histoire et l’état
enrichi ; elle ne transporte comme entrées indépendantes ni \(K\) ni \(U\).

\begin{theorem}[Carré de naturalité de la lecture signée]
\label{p12-83-c10-s1:thm:naturalidad-lector-firmado-autonomo}
On a
\begin{equation}
 \boxed{
 R_{{\rm sgn},N}\circ\tau_{N',N}^{\rm sgn}
 =\operatorname{id}_{\mathcal U_{12}^{\rm sgn}}
  \circ R_{{\rm sgn},N'}.}
 \label{p12-83-c10-s1:eq:cuadrado-naturalidad-lector-firmado}
\end{equation}
\end{theorem}

\begin{proof}
La naturalité des états conserve littéralement les cent huit premières
fenêtres du registre :
\[
 R_{12,N}\circ\tau_{N',N}^{\rm sgn}=R_{12,N'}.
\]
La troncature peut modifier uniquement l’histoire ultérieure. Comme
\(E_{108}^{90,120}\) et \(\Pi_H\) dépendent exclusivement de ces fenêtres,
les deux membres renvoient le même \(U\). La commutativité est une égalité
d’applications, et non une analogie entre diagrammes.
\end{proof}

\subsection{Orbites de pas \(3\) et transformée de Hadamard}

L’ordre cyclique du sceau se répartit dans les trois orbites
\begin{align}
 K^{(1)}&=(K_1,K_4,K_7,K_{10}),
 \label{p12-83-c10-s1:eq:orbita-k-1}\\
 K^{(2)}&=(K_2,K_5,K_8,K_{11}),
 \label{p12-83-c10-s1:eq:orbita-k-2}\\
 K^{(3)}&=(K_3,K_6,K_9,K_{12}).
 \label{p12-83-c10-s1:eq:orbita-k-3}
\end{align}
Soit
\begin{equation}
 H_4=
 \begin{pmatrix}
 1&1&1&1\\
 1&1&-1&-1\\
 1&-1&1&-1\\
 1&-1&-1&1
 \end{pmatrix}.
 \label{p12-83-c10-s1:eq:matriz-hadamard-cuatro-autonoma}
\end{equation}
Si \(P_{\rm orb}\) réordonne les douze coordonnées conformément à
\eqref{p12-83-c10-s1:eq:orbita-k-1}--\eqref{p12-83-c10-s1:eq:orbita-k-3}, la transformation globale est
\begin{equation}
 \mathcal H_{12}
 =P_{\rm orb}^{-1}
 (H_4\oplus H_4\oplus H_4)P_{\rm orb}.
 \label{p12-83-c10-s1:eq:hadamard-global-doce}
\end{equation}

Les trois blocs de \eqref{p12-83-c10-s1:eq:vector-u-firmado-autonomo}, lus par
orbites, sont
\begin{align*}
 U^{(1)}&=(2378,-452,-668,-322),\\
 U^{(2)}&=(1406,998,-204,-28),\\
 U^{(3)}&=(2479,-551,-371,-997).
\end{align*}

\begin{lemma}[Quart de tour algébrique de Hadamard]
\label{p12-83-c10-s1:lem:cuarto-inversa-hadamard}
La matrice \(H_4\) satisfait
\begin{equation}
 H_4^2=4I_4,
 \qquad
 H_4^{-1}=\frac14H_4,
 \qquad
 \det H_4=-16.
 \label{p12-83-c10-s1:eq:identidades-hadamard-cuatro}
\end{equation}
\end{lemma}

\begin{proof}
Les lignes de \(H_4\) sont orthogonales et chacune possède une norme au carré égale à quatre.
Comme la matrice est symétrique, \(H_4^{\mathsf T}H_4=H_4^2=4I_4\). La
formule de l’inverse et le déterminant s’en déduisent immédiatement.
\end{proof}

\begin{theorem}[Reconstruction entière unique du sceau]
\label{p12-83-c10-s1:thm:reconstruccion-entera-k-hadamard}
L’inversion
\begin{equation}
 K^{(r)}=\frac14H_4U^{(r)}
 \qquad(r=1,2,3)
 \label{p12-83-c10-s1:eq:inversion-bloques-hadamard}
\end{equation}
produit
\begin{align*}
 K^{(1)}&=(234,729,621,794),\\
 K^{(2)}&=(543,659,058,146),\\
 K^{(3)}&=(140,824,914,601).
\end{align*}
En défaisant l’ordre orbital, on obtient
\begin{equation}
 \boxed{
 K=(234,543,140,729,659,824,621,058,914,794,146,601).}
 \label{p12-83-c10-s1:eq:sello-k-doce-autonomo}
\end{equation}
C'est l'unique préimage rationnelle de \(U\), et elle appartient à \(\mathbb Z^{12}\).
\end{theorem}

\begin{proof}
Par \eqref{p12-83-c10-s1:eq:identidades-hadamard-cuatro}, l’antécédent rationnel est unique.
Les multiplications explicites sont
\begin{align*}
 H_4U^{(1)}&=(936,2916,2484,3176),\\
 H_4U^{(2)}&=(2172,2636,232,584),\\
 H_4U^{(3)}&=(560,3296,3656,2404).
\end{align*}
Toutes les coordonnées sont divisibles par quatre. La division et le réentrelacement
produisent \eqref{p12-83-c10-s1:eq:sello-k-doce-autonomo}.
\end{proof}

Le facteur \(1/4\) n’est pas choisi pour ajuster une sortie : c’est l’inverse
imposé par \(H_4^2=4I_4\). Ce n’est pas non plus la racine carrée de \(-1\). La
structure complexe interne \(A_W^2=-I\) apparaîtra après la construction de la
matrice de Paley ; ce sont deux faits de types différents.

\subsection{Structure entière et forme normale de Smith}

\begin{theorem}[Image entière de \(H_4\)]
\label{p12-83-c10-s1:thm:smith-hadamard-autonomo}
La forme normale de Smith est
\begin{equation}
 \operatorname{SNF}(H_4)=\operatorname{diag}(1,2,2,4).
 \label{p12-83-c10-s1:eq:snf-hadamard-autonoma}
\end{equation}
En conséquence,
\begin{equation}
 \operatorname{coker}\mathcal H_{12}
 \cong(\mathbb Z/2\mathbb Z)^6
 \oplus(\mathbb Z/4\mathbb Z)^3,
 \label{p12-83-c10-s1:eq:cociente-hadamard-doce}
\end{equation}
et l’image de \(\mathcal H_{12}\) a pour indice
\[
 16^3=4096
\]
dans \(\mathbb Z^{12}\). Pour \(u\in\mathbb Z^4\),
\begin{equation}
 u\in H_4\mathbb Z^4
 \quad\Longleftrightarrow\quad
 H_4u\equiv0\pmod4.
 \label{p12-83-c10-s1:eq:criterio-imagen-hadamard}
\end{equation}
\end{theorem}

\begin{proof}
Le plus grand commun diviseur des entrées est un. Les mineurs d’ordre deux sont
pairs et l’un d’eux a pour valeur absolue deux ; ceux d’ordre trois sont des multiples de
quatre et l’un d’eux a pour valeur absolue quatre ; le déterminant a pour module
seize. Les diviseurs déterminantaux sont \(1,2,4,16\), d’où l’on obtient
les facteurs invariants \(1,2,2,4\). La somme directe triple répète les
facteurs et multiplie les indices. Enfin,
\(H_4^{-1}=H_4/4\), de sorte que l’antécédent est entier exactement sous la
congruence \eqref{p12-83-c10-s1:eq:criterio-imagen-hadamard}.
\end{proof}

\subsection{Certification et reconstruction par différences de pas \(3\) et \(4\)}

Une fois que le registre a émis \((b_{90},b_{120},Q_{\rm TPK})\), que
\(\Pi_H\) a produit \(U\) et que Hadamard a reconstruit \(K\), les
coordonnées suivantes agissent en aval comme certificat indépendant de
ces sorties préalables. Avec des indices cycliques modulo douze, nous définissons
\begin{equation}
 (D_3K)_m=K_m-K_{m+3},
 \qquad
 (D_4K)_m=K_m-K_{m+4},
 \qquad
 Q(K)=\sum_{m=1}^{12}K_m.
 \label{p12-83-c10-s1:eq:extractor-diferencias-k}
\end{equation}
Pour \eqref{p12-83-c10-s1:eq:sello-k-doce-autonomo}, on vérifie
\begin{align*}
 D_3K=b_{90}={}&(-495,-116,-684,108,601,-90,\\
         &\qquad-173,-88,313,560,-397,461),\\
 D_4K=b_{120}={}&(-425,-281,-481,671,-255,30,\\
         &\qquad475,-543,680,251,6,-128),\\
 Q(K)=Q_{\rm TPK}={}&6263.
\end{align*}

\begin{theorem}[Complétude de l’extracteur transversal]
\label{p12-83-c10-s1:thm:completitud-extractor-transversal-k}
L'opérateur
\[
 \mathcal E_{\rm dod}=(D_3,D_4,Q):
 \mathbb Z^{12}\longrightarrow\mathbb Z^{25}
\]
a rang douze et pour forme normale de Smith
\begin{equation}
 \operatorname{SNF}(\mathcal E_{\rm dod})
 =\operatorname{diag}(\underbrace{1,\ldots,1}_{11},12).
 \label{p12-83-c10-s1:eq:snf-extractor-transversal-k}
\end{equation}
En particulier, \((D_3K,D_4K,Q(K))\) détermine un unique \(K\). Cette
unicité certifie et reconstruit les données déjà émises par le registre ; elle ne les
engendre ni ne les sélectionne.
\end{theorem}

\begin{proof}
Si \(D_3v=D_4v=0\), alors \(v\) est constant le long des pas
trois et quatre. Comme \(\gcd(3,4)=1\), ces pas engendrent tout
\(\mathbb Z/12\mathbb Z\), de sorte que
\[
 \ker D_3\cap\ker D_4=\langle\mathbf1\rangle.
\]
La charge totale ne s’annule pas sur cette droite, car \(Q(\mathbf1)=12\). Le rang
est douze. Les lignes de \((D_3,D_4)\) sont les incidences orientées d’un graphe
de Cayley connexe ; un arbre couvrant apporte onze facteurs invariants
unitaires. Tout mineur maximal doit contenir la ligne \(Q\), et les opérations
unimodulaires réduisent sa dernière contribution à douze. Il existe un mineur maximal de
module douze ; le dernier facteur est donc exactement douze.
\end{proof}

La reconstruction certificatrice de \(U\) à partir de ces coordonnées reproduit la
publication préalable de \(\Pi_H\) et ne nécessite pas de multiplier à nouveau par
\(H_4\). Pour \(r=1,2,3\), soit
\begin{equation}
 B_r=\sum_{j=0}^{3}(D_4K)_{r+3j}
 =\sum_{j=0}^{3}b_{120,r+3j}.
 \label{p12-83-c10-s1:eq:br-diferencias-k}
\end{equation}
Alors
\[
 (B_1,B_2,B_3)=(972,-1073,101).
\]
Les sommes orbitales sont
\begin{align}
 A_1&=\frac{Q+2B_1+B_2}{3}=2378,
 \label{p12-83-c10-s1:eq:a1-reconstruccion-u}\\
 A_2&=A_1-B_1=1406,
 \label{p12-83-c10-s1:eq:a2-reconstruccion-u}\\
 A_3&=A_2-B_2=2479.
 \label{p12-83-c10-s1:eq:a3-reconstruccion-u}
\end{align}
Si \(d_{r,j}=(D_3K)_{r+3j}=b_{90,r+3j}\), le bloc signé de l’orbite
\(r\) est
\begin{equation}
 \bigl(A_r,
 d_{r,1}-d_{r,3},
 d_{r,0}+d_{r,2},
 d_{r,0}-d_{r,2}\bigr),
 \label{p12-83-c10-s1:eq:bloque-firmado-desde-diferencias}
\end{equation}
ce qui reproduit les trois blocs \(U^{(r)}\) déjà publiés par \(\Pi_H\).

\subsection{Tétrade de seuil du sceau}

La complétion décimale locale satisfait
\[
 1000=729+271.
\]
Le support des coordonnées du sceau qui atteignent le seuil \(729\) est
\begin{equation}
 \boxed{
 B_K:=\{m:K_m\geq729\}=\{4,6,9,10\}.}
 \label{p12-83-c10-s1:eq:tetrada-umbral-k}
\end{equation}
Cette tétrade s’obtient après la construction de \(K\). Dans cette unité, elle n’est pas
encore appelée étoile de Witt : un tel objet ne sera défini qu’après
la construction du code ternaire, de ses hexades et du système de Steiner.

\subsection{Indépendance de la construction généalogique à l'égard de données métrologiques ou physiques externes}

\begin{theorem}[Isolement de \(U\) et \(K\)]
\label{p12-83-c10-s1:thm:aislamiento-u-k-alpha-codata}
La composition
\begin{equation}
 \Omega_{\rm seed}\longrightarrow x_{\rm term}
 \xrightarrow{R_{\rm sgn}}U
 \xrightarrow{\mathcal H_{12}^{-1}}K
 \label{p12-83-c10-s1:eq:composicion-upstream-u-k}
\end{equation}
ne reçoit en entrée ni \(\alpha\), ni \(\alpha^{-1}\), ni CODATA, ni une suite
décimale cible. Le sceau \(K\) est déterminé avant la définition de la
récurrence qui produira \(\alpha\).
\end{theorem}

\begin{proof}
La première partie de la composition utilise uniquement des germes, des émissions,
des lecteurs ternaires, des opérateurs \(L_0,L_1\), des relations de survie,
le transport de phase et les cartes de l’état TPK. La lecture
\(R_{\rm sgn}=\Pi_H\circ E_{108}^{90,120}\circ R_{12}\) est la composition
causale du registre terminal, et la dernière flèche est la transformation linéaire
\(\mathcal H_{12}^{-1}=\mathcal H_{12}/4\)
sur son image entière. Aucune des formules ne contient une variable
physique ou un tableau métrologique. La récurrence en base mille n’est introduite que
dans l’unité suivante, après avoir fixé \(K\).
\end{proof}

\subsection{Obstructions du registre dodécaphasique et du sceau \(K\) : naturalité, intégrité de Hadamard et isolement métrologique}

\begin{criterion}[Réfutation du registre et du sceau]
La construction est réfutée si l'un des faits
suivants est vérifié :
\begin{enumerate}
 \item les fenêtres \(E_m\) ne partitionnent pas les cent huit pas ;
 \item la carte \(\mathcal L_{12}\) ne respecte pas son inverse ou ses congruences ;
 \item le recensement terminal ne produit pas \(331\) triplets, deux matrices et une unique
       charge axiale ;
 \item on tente d’obtenir \(U\) à partir de \(B_{\rm term}\) après avoir oublié les
       coordonnées signées ;
 \item le carré de naturalité
       \eqref{p12-83-c10-s1:eq:cuadrado-naturalidad-lector-firmado} échoue ;
 \item \(H_4^2\neq4I_4\), un antécédent n’est pas entier ou le vecteur
       reconstruit diffère de \eqref{p12-83-c10-s1:eq:sello-k-doce-autonomo} ;
 \item la forme normale de Smith diffère de
       \eqref{p12-83-c10-s1:eq:snf-hadamard-autonoma} ;
 \item le certificat \((D_3,D_4,Q)\) perd du rang, ne coïncide pas avec les
       émissions préalables du registre ou ne reproduit pas \(U\) ;
 \item \(B_K\) est introduit avant la construction de \(K\), ou on lui attribue une
       structure de Witt avant la définition de \(W_{12}\) ;
 \item un chiffre de \(\alpha\) ou une référence métrologique intervient dans
       l’une quelconque des flèches de \eqref{p12-83-c10-s1:eq:composicion-upstream-u-k}.
\end{enumerate}
\end{criterion}


% Unidad U017. Cierre afín, prolongación y semántica de frontera de alpha.
% Redacción autónoma desde la autoridad material 1063 y su sucesor V2.

\section{Construction initiale de la frontière \texorpdfstring{\(662/663\)}{662/663} au moyen du cocycle dodécaphasique de retenue}
\label{p12-83-c10-s2:chap:acarreo-alpha-frontera-662-663}

\subsection{Domaine arithmétique de la clôture}

La construction reçoit quatre coordonnées dodécaphasiques obtenues dans les
unités précédentes :
\[
\begin{aligned}
 P={}&(141,592,653,589,793,238,462,643,383,279,502,884),\\
 E={}&(718,281,828,459,045,235,360,287,471,352,662,497),\\
 \Phi={}&(618,033,988,749,894,848,204,586,834,365,638,117),\\
 K={}&(234,543,140,729,659,824,621,058,914,794,146,601).
\end{aligned}
\]
Les trois premiers vecteurs sont les douze premières triades des lectures
archimédiennes fractionnaires de \(\pi\), \(e\) et \(\varphi\). Le quatrième est
la coordonnée dodécaphasique de l’état terminal enrichi du Topological
Phase Kernel, construite avant l’ouverture de l’opération de retenue.

Nous fixons la base
\[
 B:=10^3=1000
\]
et la concaténation
\begin{equation}
 \iota(v_1,\ldots,v_{12})
 :=\sum_{m=1}^{12}v_mB^{12-m}.
 \label{p12-83-c10-s2:eq:concatenacion-base-mil-alpha}
\end{equation}
\Needspace{8\baselineskip}
En particulier,
\[
\begin{aligned}
 \iota(P)&=141592653589793238462643383279502884,\\
 \iota(E)&=718281828459045235360287471352662497,\\
 \iota(\Phi)&=618033988749894848204586834365638117,\\
 \iota(K)&=234543140729659824621058914794146601.
\end{aligned}
\]

L’opération étudiée dans cette unité est l’application
\begin{equation}
 \mathfrak C_{12}:
 (P,E,\Phi,K,c_{13})
 \longmapsto(A,C),
 \label{p12-83-c10-s2:eq:mapa-cierre-dodecafasico-alpha}
\end{equation}
où \(A=(A_1,\ldots,A_{12})\) est le mot de sortie et
\[
 C=(c_1,\ldots,c_{13})
\]
est le système complet des treize quotients euclidiens : douze quotients
internes et une condition à la frontière droite.

\subsection{Division euclidienne orientée de droite à gauche}

Pour \(m=12,11,\ldots,1\), on définit
\begin{equation}
\begin{aligned}
 s_m&:=P_m+E_m-\Phi_m-K_m+c_{m+1},\\
 A_m&:=s_m\bmod B,
       \qquad 0\leq A_m<B,\\
 c_m&:=\frac{s_m-A_m}{B}
      =\left\lfloor\frac{s_m}{B}\right\rfloor .
\end{aligned}
\label{p12-83-c10-s2:eq:recurrencia-corta-alpha}
\end{equation}
La deuxième égalité pour \(c_m\) est également valable lorsque \(s_m<0\),
car on utilise la division euclidienne avec reste dans
\(\{0,\ldots,B-1\}\). C'est pourquoi,
\[
 -431=569-1000,
 \qquad
 -717=283-1000,
 \qquad
 -1200=800-2\cdot1000,
\]
et les emprunts négatifs n’introduisent aucune ambiguïté.

\begin{theorem}[Existence et unicité de la retenue orientée]
\label{p12-83-c10-s2:thm:unicidad-acarreo-orientado-alpha}
Une fois \(P,E,\Phi,K\) et la retenue terminale \(c_{13}\) fixés, la récurrence
\eqref{p12-83-c10-s2:eq:recurrencia-corta-alpha} détermine de manière unique
\[
 (A_1,\ldots,A_{12};c_1,\ldots,c_{12}).
\]
\end{theorem}

\begin{proof}
La ligne \(m=12\) est déterminée par la division euclidienne de
\[
 P_{12}+E_{12}-\Phi_{12}-K_{12}+c_{13}
\]
par \(B\). Une fois \(c_{12}\) connu, la ligne \(m=11\) est déterminée par
la même propriété. La récurrence descendante atteint \(m=1\). À chaque
étape, l’unicité du quotient et du reste euclidiens empêche deux sorties
distinctes.
\end{proof}

\subsection{Fenêtre dodécaphasique de l'état signé et clôture réversible du cycle}

La fenêtre finie est définie par la condition terminale
\[
 c_{13}^{\rm fin}=0.
\]
L’évaluation complète est la suivante. Le tableau est imprimé dans l’ordre
croissant, bien qu’il doive être lu causalement de la dernière ligne vers la
première.

\begingroup
\scriptsize
\setlength{\tabcolsep}{3.2pt}
\begin{longtable}{@{}r|rrrr|r|r|r|r@{}}
\caption{Les douze divisions euclidiennes et les treize retenues de la fenêtre
dodécaphasique finie.}
\label{p12-83-c10-s2:tab:acarreo-alpha-finito-autonomo}\\
\toprule
\(m\)&\(P_m\)&\(E_m\)&\(\Phi_m\)&\(K_m\)&
\(c_{m+1}\)&\(s_m\)&\(A_m\)&\(c_m\)\\
\midrule
\endfirsthead
\toprule
\(m\)&\(P_m\)&\(E_m\)&\(\Phi_m\)&\(K_m\)&
\(c_{m+1}\)&\(s_m\)&\(A_m\)&\(c_m\)\\
\midrule
\endhead
1  &141&718&618&234& \(0\)& \(7\)&007& \(0\)\\
2  &592&281&033&543& \(0\)& \(297\)&297& \(0\)\\
3  &653&828&988&140& \(-1\)& \(352\)&352& \(0\)\\
4  &589&459&749&729& \(-1\)& \(-431\)&569& \(-1\)\\
5  &793&045&894&659& \(-2\)& \(-717\)&283& \(-1\)\\
6  &238&235&848&824& \(-1\)& \(-1200\)&800& \(-2\)\\
7  &462&360&204&621& \(0\)& \(-3\)&997& \(-1\)\\
8  &643&287&586&058& \(-1\)& \(285\)&285& \(0\)\\
9  &383&471&834&914& \(-1\)& \(-895\)&105& \(-1\)\\
10 &279&352&365&794& \(0\)& \(-528\)&472& \(-1\)\\
11 &502&662&638&146& \(0\)& \(380\)&380& \(0\)\\
12 &884&497&117&601& \(0\)& \(663\)&663& \(0\)\\
\bottomrule
\end{longtable}
\endgroup

Le mot complet des treize retenues est
\[
 \boxed{
 C^{\rm fin}
 =(c_1,\ldots,c_{13})
 =(0,0,0,-1,-1,-2,-1,0,-1,-1,0,0,0).}
\]
La sortie dodécaphasique est
\[
 \boxed{
 A^{\rm fin}
 =(007,297,352,569,283,800,997,285,105,472,380,663).}
\]
Nous définissons le rationnel fini
\begin{equation}
 \boxed{
 \alpha_{12}^{\rm fin}
 :=10^{-36}\iota(A^{\rm fin})
 =\frac{7297352569283800997285105472380663}{10^{36}}.}
 \label{p12-83-c10-s2:eq:alpha-finita-autonoma}
\end{equation}
Par conséquent,
\[
 \alpha_{12}^{\rm fin}
 =0.007297352569283800997285105472380663.
\]

\subsection{Reconnaissance métrologique ultérieure : CODATA 2018}

La comparaison métrologique n’est effectuée qu’après clôture de la récurrence,
de la frontière droite et du rationnel précédent. L’identité entière donne
\begin{equation}
 \boxed{
 \alpha_{12}^{\rm fin}
 =10^{-36}
 \left\lfloor\frac{10^{36}}{137.035999084}\right\rfloor.}
 \label{p12-83-c10-s2:eq:alpha-finita-codata-2018}
\end{equation}
En conséquence,
\begin{equation}
 (\alpha_{12}^{\rm fin})^{-1}
 =137.03599908400000000000000000000001066588\ldots,
 \label{p12-83-c10-s2:eq:alpha-inversa-codata-2018}
\end{equation}
et reproduit exactement la suite centrale publiée par CODATA 2018
\cite{Tiesinga2021CODATA} :
\[
 \boxed{137.035999084.}
\]
Le mot \emph{exactement} renvoie à la suite centrale publiée, et non
à une prétendue mesure expérimentale à trente-six chiffres. Les chiffres
ultérieurs de \eqref{p12-83-c10-s2:eq:alpha-inversa-codata-2018} sont des conséquences du
rationnel HMT déjà fixé. CODATA ne sélectionne pas l’état, ne détermine pas (K),
ne décide pas du sens de la retenue et n’intervient pas dans le tableau dodécaphasique.

\subsection{Identité affine complète}

Chaque ligne du tableau satisfait exactement
\begin{equation}
 P_m+E_m-\Phi_m-K_m+c_{m+1}
 =A_m+Bc_m.
 \label{p12-83-c10-s2:eq:identidad-local-acarreo-alpha}
\end{equation}
En introduisant
\[
 C^-:=(c_1,\ldots,c_{12}),
 \qquad
 C^+:=(c_2,\ldots,c_{13}),
\]
les douze identités locales se réunissent en
\begin{equation}
 \boxed{
 P+E-\Phi-K+C^+-BC^-=A.}
 \label{p12-83-c10-s2:eq:identidad-afin-vectorial-alpha}
\end{equation}
Le terme
\[
 \partial_BC:=C^+-BC^-
\]
est le cocycle de retenue de la chaîne orientée. Par conséquent,
\(P+E-\Phi-K=A\) n’est pas une égalité coordonnée par coordonnée dans
\(\mathbb Z^{12}\) : elle ne devient correcte qu’en incorporant
\(\partial_BC\).

\begin{theorem}[Télescopage des treize retenues]
\label{p12-83-c10-s2:thm:telescopado-acarreos-alpha}
Pour toute condition terminale \(c_{13}\), la récurrence satisfait
\begin{equation}
 \boxed{
 \iota(P)+\iota(E)-\iota(\Phi)-\iota(K)-\iota(A)
 =B^{12}c_1-c_{13}.}
 \label{p12-83-c10-s2:eq:telescopado-general-alpha}
\end{equation}
Dans la fenêtre finie, \(c_1=c_{13}=0\), et donc
\begin{equation}
 \boxed{
 \iota(P)+\iota(E)-\iota(\Phi)-\iota(K)
 =\iota(A^{\rm fin}).}
 \label{p12-83-c10-s2:eq:identidad-entera-alpha-finita}
\end{equation}
\end{theorem}

\begin{proof}
Nous multiplions \eqref{p12-83-c10-s2:eq:identidad-local-acarreo-alpha} par
\(B^{12-m}\) et sommons sur \(m=1,\ldots,12\). La contribution des
retenues est
\[
 \sum_{m=1}^{12}
 (Bc_m-c_{m+1})B^{12-m}.
\]
Les composantes intérieures s’annulent :
\begin{align*}
 \sum_{m=1}^{12}(Bc_m-c_{m+1})B^{12-m}
 &=\sum_{m=1}^{12}c_mB^{13-m}
   -\sum_{m=1}^{12}c_{m+1}B^{12-m}\\
 &=B^{12}c_1-c_{13}.
\end{align*}
Cela prouve \eqref{p12-83-c10-s2:eq:telescopado-general-alpha}. La spécialisation
\(c_1=c_{13}=0\) donne \eqref{p12-83-c10-s2:eq:identidad-entera-alpha-finita}.
\end{proof}

\subsection{Inversion exacte du sceau \texorpdfstring{\(K\)}{K}}

La récurrence locale est réversible. En isolant \(K_m\) dans
\eqref{p12-83-c10-s2:eq:identidad-local-acarreo-alpha},
\begin{equation}
 \boxed{
 K_m=P_m+E_m-\Phi_m+c_{m+1}-A_m-Bc_m.}
 \label{p12-83-c10-s2:eq:inversion-local-k-desde-alpha}
\end{equation}
Par concaténation,
\begin{equation}
 \boxed{
 \iota(K)
 =\iota(P)+\iota(E)-\iota(\Phi)-\iota(A)
  -B^{12}c_1+c_{13}.}
 \label{p12-83-c10-s2:eq:inversion-concatenada-k-desde-alpha}
\end{equation}
À la frontière finie,
\[
 \iota(K)
 =\iota(P)+\iota(E)-\iota(\Phi)-\iota(A^{\rm fin}).
\]

La réversibilité algébrique n’inverse pas l’ordre causal. L’ordre effectif
est
\[
 x_{\rm term}
 \longrightarrow K
 \longrightarrow(P,E,\Phi,K,c_{13})
 \longrightarrow(A,C).
\]
La formule \eqref{p12-83-c10-s2:eq:inversion-local-k-desde-alpha} est un contrôle de
commutativité : elle démontre que la sortie et les retenues récupèrent le sceau
qui était déjà fixé. Elle ne constitue pas une règle permettant de sélectionner
rétrospectivement \(K\) à partir d’une valeur souhaitée de \(\alpha\).

\subsection{Prolongement coinductif du sceau dodécaphasique et récurrence inverse de la retenue dans \(\mathcal C=\{-2,-1,0,1,2\}\)}

Soient
\[
 (P_k)_{k\geq1},
 \qquad
 (E_k)_{k\geq1},
 \qquad
 (\Phi_k)_{k\geq1}
\]
les suites infinies de triades produites par les trois branches
coinductives, et soit
\[
 \kappa(k):=1+((k-1)\bmod12).
\]
Le prolongement du sceau est périodique :
\[
 K_k^{\rm per}:=K_{\kappa(k)}.
\]
Pour une profondeur finie \(N\) et une condition terminale
\(t=c_{N+1}\), la récurrence est
\begin{equation}
\begin{aligned}
 s_k^{(N,t)}
 &=P_k+E_k-\Phi_k-K_{\kappa(k)}+c_{k+1}^{(N,t)},\\
 A_k^{(N,t)}
 &=s_k^{(N,t)}\bmod1000,\\
 c_k^{(N,t)}
 &=\left\lfloor\frac{s_k^{(N,t)}}{1000}\right\rfloor,
 \qquad k=N,\ldots,1.
\end{aligned}
\label{p12-83-c10-s2:eq:recurrencia-larga-alpha}
\end{equation}

L’ensemble des retenues terminales
\[
 \mathcal C:=\{-2,-1,0,1,2\}
\]
est invariant pour les entrées certifiées : pour chaque triade et chaque
\(c_{k+1}\in\mathcal C\), le quotient \(c_k\) appartient à nouveau à
\(\mathcal C\). On ne choisit pas une condition terminale favorable. On propage
les cinq conditions et l’on ne publie que le préfixe commun à toutes.

L’exécution à profondeur certifiée emploie
\[
\begin{array}{c|r}
\text{quantité}&\text{valeur}\\ \hline
\text{blocs ternaires par canal}&3501\\
\text{trits par canal}&3501\cdot6=21006\\
\text{chiffres décimaux de garde}&10020\\
\text{triades de garde}&3340\\
\text{conditions terminales propagées}&5\\
\text{triades communes}&3339\\
\text{chiffres communs}&10017\\
\text{chiffres publiés}&10000
\end{array}
\]
Les cinq branches terminales coalescent avant le préfixe publié. En
particulier, aucune condition terminale extraite d’une valeur cible de
\(\alpha\) n’intervient dans la sortie.

\begin{theorem}[Compatibilité des sorties stabilisées]
\label{p12-83-c10-s2:thm:compatibilidad-prolongacion-alpha}
Soient \(A^{[N]}\) et \(A^{[N']}\), avec \(N'>N\), deux exécutions dont les
cinq conditions terminales ont coalescé avant la position \(N\).
Alors la troncature de \(A^{[N']}\) à ses \(N\) premiers blocs
coïncide avec \(A^{[N]}\). Les préfixes stabilisés définissent une unique
section dans la limite inverse des mots finis.
\end{theorem}

\begin{proof}
La récurrence est déterministe de droite à gauche une fois connue la
retenue entrant dans une position. La coalescence affirme que toutes les
conditions distantes produisent la même retenue avant la frontière du
préfixe. À partir de ce point, les mêmes triades \(P_k,E_k,\Phi_k\) et le même
sceau périodique produisent, par unicité euclidienne, les mêmes restes et
quotients à toutes les positions précédentes. C’est exactement la
commutativité avec la troncature.
\end{proof}

\subsection{Frontière provenant de la queue distante}

Le prolongement détermine à la frontière visible
\[
 c_{13}^{\infty}=-1.
\]
Les douze retenues intérieures restent identiques à celles de la fenêtre
finie :
\[
 \boxed{
 C_{\leq13}^{\infty}
 =(0,0,0,-1,-1,-2,-1,0,-1,-1,0,0,-1).}
\]
Seule change l’information qui entre depuis la treizième triade. Les
deux divisions pertinentes sont
\[
\begin{aligned}
 s_{12}^{\infty}
 &=884+497-117-601-1=662,\\
 A_{12}^{\infty}&=662,
 \qquad c_{12}^{\infty}=0,
\end{aligned}
\]
et
\[
\begin{aligned}
 s_{13}^{\infty}
 &=197+757-720-234-1=-1,\\
 A_{13}^{\infty}&=999,
 \qquad c_{13}^{\infty}=-1.
\end{aligned}
\]
Par conséquent, le début du mot prolongé est
\[
 \boxed{
 A^{\infty}
 =(007,297,352,569,283,800,997,285,105,472,380,662,999,\ldots).}
\]
Le développement correspondant commence par
\begin{equation}
 \boxed{
 \alpha_{\infty}
 =0.007297352569283800997285105472380662999564172539642\ldots.}
 \label{p12-83-c10-s2:eq:alpha-infinita-autonoma}
\end{equation}

La différence entre les deux frontières est locale et exacte :
\[
 c_{13}^{\rm fin}-c_{13}^{\infty}=1,
 \qquad
 A_{12}^{\rm fin}-A_{12}^{\infty}=1,
\]
tandis que
\[
 A_m^{\rm fin}=A_m^\infty,
 \qquad 1\leq m\leq11.
\]

En appliquant le théorème de télescopage aux douze premiers blocs du
prolongement,
\[
 \iota(P)+\iota(E)-\iota(\Phi)-\iota(K)
 -\iota(A^\infty_{\leq12})
 =-c_{13}^{\infty}=1.
\]
Donc
\begin{equation}
 \boxed{
 \iota(A^{\rm fin})
 =\iota(A^\infty_{\leq12})+1.}
 \label{p12-83-c10-s2:eq:relacion-unidad-fronteras-alpha}
\end{equation}
La terminaison \(663\) ne contredit pas la terminaison \(662\) : la première
est le représentant fini obtenu en fermant la fenêtre avec
\(c_{13}=0\) ; la seconde conserve la retenue \(-1\) provenant de la queue
distante.

\begin{figure}[htbp]
\centering
\begin{tikzpicture}[
  >=Stealth,
  node distance=12mm and 16mm,
  box/.style={draw,rounded corners=1.5mm,align=center,
    inner xsep=3.2mm,inner ysep=2.5mm},
  arr/.style={->,thick}
]
\node[box] (shared) {données partagées\\
 \(P,E,\Phi,K\)};
\node[box,below left=of shared] (finite)
 {frontière finie\\\(c_{13}=0\)};
\node[box,below right=of shared] (remote)
 {queue distante\\\(c_{13}=-1\)};
\node[box,below=of finite] (a663)
 {\(A_{12}=663\)\\clausura finita};
\node[box,below=of remote] (a662)
 {\(A_{12}=662\), \(A_{13}=999\)\\préfixe prolongé};
\draw[arr] (shared)--(finite);
\draw[arr] (shared)--(remote);
\draw[arr] (finite)--(a663);
\draw[arr] (remote)--(a662);
\draw[<->,densely dashed] (a663)--node[below,align=center]
 {différence entière exacte \(1\)} (a662);
\end{tikzpicture}
\caption{Bifurcation de frontière d’une même récurrence. On ne compare pas
deux générateurs distincts : seule change la donnée qui entre depuis la
droite.}
\label{p12-83-c10-s2:fig:bifurcacion-frontera-alpha}
\end{figure}

\subsection{Troncature et arrondi à \(36\) chiffres}

Pour \(x\geq0\), nous définissons
\[
 \operatorname{Tr}_{36}(x)
 :=10^{-36}\left\lfloor10^{36}x\right\rfloor
\]
et
\[
 \operatorname{Rd}_{36}(x)
 :=10^{-36}\left\lfloor10^{36}x+\frac12\right\rfloor.
\]
Soit
\[
 a:=7297352569283800997285105472380662.
\]
Le développement \eqref{p12-83-c10-s2:eq:alpha-infinita-autonoma} peut s’écrire
\[
 \alpha_\infty=\frac{a+\rho}{10^{36}},
 \qquad
 \rho=0.999564172539642709\ldots,
\]
d’où \(0.999<\rho<1\).

\begin{theorem}[Résolution exacte de la frontière \(662/663\)]
\label{p12-83-c10-s2:thm:frontera-alpha-662-663}
Le prolongement coinductif satisfait
\[
 \boxed{
 \operatorname{Tr}_{36}(\alpha_\infty)
 =0.007297352569283800997285105472380662,}
\]
tandis que son arrondi à trente-six chiffres satisfait
\[
 \boxed{
 \operatorname{Rd}_{36}(\alpha_\infty)
 =0.007297352569283800997285105472380663
 =\alpha_{12}^{\rm fin}.}
\]
\end{theorem}

\begin{proof}
Comme \(0\leq\rho<1\),
\[
 \left\lfloor10^{36}\alpha_\infty\right\rfloor=a.
\]
Cela donne la troncature se terminant par \(662\). Comme \(\rho>1/2\),
\[
 \left\lfloor10^{36}\alpha_\infty+\frac12\right\rfloor=a+1,
\]
qui se termine par \(663\). Par \eqref{p12-83-c10-s2:eq:alpha-finita-autonoma},
\((a+1)10^{-36}\) est exactement \(\alpha_{12}^{\rm fin}\).
\end{proof}

De plus,
\[
 0<\alpha_{12}^{\rm fin}-\alpha_\infty
 =\frac{1-\rho}{10^{36}}<10^{-39}.
\]
La différence n’est pas une erreur arithmétique du tableau : elle quantifie la
différence entre deux opérateurs de frontière distincts.

\subsection{Isolement causal à l’égard de \texorpdfstring{\(\alpha\)}{alpha} et des données métrologiques}

Pour la profondeur \(N\), le domaine complet de l’opération est
\[
 \mathscr D_N
 =\bigl((P_k)_{k=1}^{N},(E_k)_{k=1}^{N},
        (\Phi_k)_{k=1}^{N},K,c_{N+1}\bigr).
\]
L’application de retenue
\[
 \mathfrak C_N:\mathscr D_N
 \longrightarrow
 \{0,\ldots,999\}^{N}\times\mathcal C^{N}
\]
est définie exclusivement par \eqref{p12-83-c10-s2:eq:recurrencia-larga-alpha}. Son
domaine ne contient ni valeur cible de \(\alpha\), ni son inverse, ni tableau
métrologique, ni suite décimale cible, ni retenue terminale choisie par
comparaison.

\begin{theorem}[Non-interférence de la cible métrologique]
\label{p12-83-c10-s2:thm:no-interferencia-metrologica-alpha}
Soit \(\mathscr Y\) un espace quelconque de données externes contenant des valeurs
métrologiques ou des tableaux de comparaison. L’extension formelle
\[
 \widetilde{\mathfrak C}_N:
 \mathscr D_N\times\mathscr Y
 \longrightarrow
 \{0,\ldots,999\}^{N}\times\mathcal C^{N},
 \qquad
 \widetilde{\mathfrak C}_N(d,y):=\mathfrak C_N(d),
\]
se factorise par la projection
\[
 \operatorname{pr}_{\mathscr D_N}:
 \mathscr D_N\times\mathscr Y\longrightarrow\mathscr D_N.
\]
Par conséquent, pour tous \(y,y'\in\mathscr Y\),
\[
 \widetilde{\mathfrak C}_N(d,y)
 =\widetilde{\mathfrak C}_N(d,y').
\]
Aucune intervention sur une valeur de \(\alpha\) ou sur un tableau
métrologique n’altère \(K\), les quotients, les restes ni le préfixe produit.
\end{theorem}

\begin{proof}
L’affirmation résulte de
\[
 \widetilde{\mathfrak C}_N
 =\mathfrak C_N\circ\operatorname{pr}_{\mathscr D_N}.
\]
Toutes les opérations de \(\mathfrak C_N\) sont des additions, des soustractions et des divisions
euclidiennes des éléments de \(\mathscr D_N\). Aucune coordonnée de
\(\mathscr Y\) n’y apparaît.
\end{proof}

L’ordre causal complet est
\[
 \text{branches enrichies de }\pi,e,\varphi
 \longrightarrow(P,E,\Phi),
\]
\[
 x_{\rm term}\longrightarrow K,
\]
\[
 (P,E,\Phi,K,\text{frontière})
 \longrightarrow(A,C)
 \longrightarrow\alpha_\infty,
\]
\[
 \alpha_\infty
 \longrightarrow\text{comparaison externe}.
\]
La dernière flèche ne peut pas être inversée pour sélectionner l’une quelconque des
précédentes.

\subsection{Obstructions de la prolongation affine de \(\alpha\) : télescopage, inversion de \(K\), frontière \(662/663\) et isolement des données cibles}

\begin{criterion}[Réfutation de la clôture affine et de sa prolongation]
Le bloc est réfuté dans son domaine si l'un des faits
suivants se produit :
\begin{enumerate}
 \item une ligne de la \cref{p12-83-c10-s2:tab:acarreo-alpha-finito-autonomo}
       ne satisfait pas \(s_m=A_m+1000c_m\) ;
 \item le vecteur publié des treize retenues ne reproduit pas les douze lignes ;
 \item l’identité télescopique échoue
       \[
       \iota(P)+\iota(E)-\iota(\Phi)-\iota(K)-\iota(A)
       =1000^{12}c_1-c_{13};
       \]
 \item l’inversion \eqref{p12-83-c10-s2:eq:inversion-local-k-desde-alpha} ne récupère pas
       les douze coordonnées de \(K\) ;
 \item la propagation distante produit \(c_{13}\neq-1\) pour le préfixe
       certifié ;
 \item les cinq conditions terminales de
       \(\mathcal C=\{-2,-1,0,1,2\}\) ne coalescent pas avant les chiffres
       publiés ;
 \item la triade suivant \(662\) n’est pas \(999\), auquel cas l’identification
       de l’arrondi à \(663\) cesse d’être démontrée ;
 \item le générateur ouvre un développement cible de \(\alpha\), son inverse,
       un tableau métrologique ou une condition terminale choisie pour imposer
       la coïncidence.
\end{enumerate}
\end{criterion}

Deux identités de frontière distinctes sont ainsi démontrées :
\[
 \boxed{
 663=\text{clôture finie avec }c_{13}=0
     =\text{arrondi à 36 chiffres de }\alpha_\infty,}
\]
\[
 \boxed{
 662=\text{troncature à 36 chiffres de }\alpha_\infty
 \quad\text{avec}\quad c_{13}=-1,\ A_{13}=999.}
\]
Le résultat métrologique central est finalement publié dans la même
résidence que la dérivation :
\begin{equation}
 \boxed{
 \begin{gathered}
  (\alpha_{12}^{\rm fin})^{-1}
  =137.03599908400000000000000000000001066588\ldots,\\
  \text{chaîne centrale CODATA 2018 :}\qquad137.035999084.
 \end{gathered}}
 \label{p12-83-c10-s2:eq:resultado-final-alpha-codata-2018}
\end{equation}
L’implication exprime une comparaison ultérieure avec la suite centrale
publiée ; elle n’introduit pas cette suite dans le constructeur.
